\section{来源审计：这是一张课程地图，不是一份模型年表}

本讲是 CS25 Transformers United V1 的导论，由 Advay Pal、Chetanya Rastogi 与 Div Garg 主讲。课程在 2021 年秋季开设，公开视频于 2022 年 7 月上传。原视频只有 22 分 43 秒，却承担三项任务：说明 speaker series 的组织方式，给出 attention 到 Transformer 的最小历史线索，并用 GPT-3 与 BERT 建立 decoder-only/encoder-only 的家族坐标。

仓库没有找到独立 slide PDF，但视频清楚展示了原始 deck。本文从官方视频恢复 17 张有教学信息的 slide；纯章节分隔页与 self-attention 动画的重复 build state 被有意省略。扩展公式、手算例子与伪代码用于让讲义可独立学习，不表示每一行都在课堂上逐字出现。

\begin{warningbox}{历史快照与当前知识}
“Where we are”与“The Future”反映的是 2021 年课堂视角。本文保留它们作为 field snapshot，不把当时的未来判断写成 2026 年当前事实。后续技术扩展只用于解释持久机制，并与课堂原话分开。
\end{warningbox}

\subsection{Speaker series 的学习接口}

CS25 V1 不是按照一部教材逐章讲授，而是邀请不同领域研究者说明 Transformer 如何进入自己的问题。导论的作用因此类似 API contract：先给出共同术语与架构，再让后续讲者分别替换数据、目标、mask、tokenization 与 evaluation。读者需要持续追问“哪些结构是 Transformer 的，哪些是领域特定设计”。

\lecturefigure{slide-01-course-page.jpg}{CS25 V1 以跨领域 guest lectures 组织 Transformer 知识。}{官方视频 00:30；Stanford CS25 V1 course page。}

\begin{knowledgebox}{读图：课程页不是行政截图，而是知识拓扑}
页面把 instructors、faculty advisor 与 speaker schedule 放在一起，说明课程的主要证据来自一线研究实践。后续讲次可能覆盖 NLP、vision、reinforcement learning、generative modeling 或 biology，但共同语言仍是 token 表示、attention pattern、训练目标和系统代价。阅读整门课时，先记录每位讲者改变了哪个接口，再比较哪些机制跨领域复用。
\end{knowledgebox}

\lecturefigure{slide-02-title.jpg}{Transformers United 的核心对象是深度学习模型，而不是会变形的机器人。}{官方视频 00:55。}

\begin{knowledgebox}{读图：玩笑背后的范围声明}
标题用 Transformers 玩具形象制造记忆点，但讲者立即把范围收回到 deep learning。真正的“变形”是计算图的可迁移性：同一个 token-mixing 核心可以处理单词、图像 patch、动作、音频 frame 或生物序列；输入语义不同，表示与目标也不同。课程不会把所有领域简化成同一任务，而是研究统一骨架如何与领域归纳偏置结合。
\end{knowledgebox}

\teachervoice{开场特意澄清“不是会变成汽车的机器人”，随后强调 Transformer 从 NLP 扩散到 CV 与 RL。这个轻松过渡保存了课程的真正 thesis：模型名相同不等于任务相同，但可复用的 interaction mechanism 正在连接原本分散的研究社区。}

\subsection{谁在讲，以及为什么跨领域}

三位讲师的背景覆盖 production software、generative modeling、reinforcement learning、robotics 与 NLP。这种组合影响了课程的选题：既关心公式，也关心模型如何进入真实系统；既讨论语言，也鼓励把 architecture 看成可迁移的研究工具。

\lecturefigure{slide-03-instructors.jpg}{三位讲师的背景横跨软件工程、NLP、生成建模、RL 与 robotics。}{官方视频 01:30。}

\begin{knowledgebox}{读图：背景差异决定问题差异}
软件工程视角会追问部署、接口与可维护性；NLP 视角关心语言表示、预训练和迁移；RL/robotics 视角关心时序决策、部分可观测与行动反馈。Transformer 并不会自动解决这些问题，但它提供统一的 sequence interaction primitive，使不同领域能够共享优化、规模化和表示学习经验。
\end{knowledgebox}

\subsection{三项学习目标}

导论明确给出三个 outcome：理解 Transformer 如何工作；理解它为何超越 NLP；从 guest talks 中发现新研究方向。它们分别对应 mechanism、transfer 与 research taste。只记住 $QK^\top V$ 不能完成第二、第三项；只追逐应用名词又无法判断迁移是否合理。

\lecturefigure{slide-04-learning-goals.jpg}{课程目标依次是机制理解、跨领域应用与新研究方向。}{官方视频 02:05。}

\begin{knowledgebox}{读图：把目标变成三层检查表}
Mechanism 层要求解释 token 如何交互、mask 如何限制信息、position 如何进入表示；application 层要求说明 tokenization、objective、data 与 evaluation 的变化；research 层要求识别当前失败模式和可证伪假设。后续每讲都可用这三层复盘，避免把 guest lecture 变成公司或论文名录。
\end{knowledgebox}

\subsection{本章小结}

Lecture 01 的角色是建立共同接口。课程用 guest talks 连接多个领域；本文则用完整来源幻灯片、教师口吻和自包含推导把这张地图补全。

